\documentclass[11pt]{article}
\pagestyle{plain}

% \pgfplotsset{compat=1.18}
\usepackage[margin=1in]{geometry}
\usepackage{amsmath,amsfonts,amssymb,amsthm}
\usepackage[ruled,vlined,linesnumbered,commentsnumbered]{algorithm2e}
\SetKwInput{KwIn}{Input}
\SetKwInput{KwOut}{Output}
\SetKwInput{KwGlobal}{Global variable}
\SetKwInput{KwOracle}{Oracle}
\SetKwComment{Comment}{// }{}
\usepackage{graphicx}
\usepackage{enumerate}
\usepackage{bbm}
\usepackage{verbatim}
\usepackage{hyperref,color}
\usepackage[capitalize,nameinlink]{cleveref}
\usepackage[dvipsnames]{xcolor}
\hypersetup{
	colorlinks=true,
	pdfpagemode=UseNone,
	citecolor=OliveGreen,
	linkcolor=NavyBlue,
	urlcolor=Magenta,
	pdfstartview=FitW
}
\usepackage{appendix}
\usepackage{makecell}
\usepackage{tablefootnote}
\crefname{appsec}{Appendix}{Appendices}
\usepackage{tikz}
\usepackage{pgfplots}
\usepackage{xifthen}
\usepackage{subcaption}
\usepackage{hhline}

\theoremstyle{plain}
\newtheorem{theorem}{Theorem}[section]
\newtheorem{proposition}[theorem]{Proposition}
\newtheorem{lemma}[theorem]{Lemma}
\newtheorem{corollary}[theorem]{Corollary}
\newtheorem{conjecture}[theorem]{Conjecture}
\newtheorem{problem}[theorem]{Problem}
\newtheorem{claim}[theorem]{Claim}
\newtheorem{fact}[theorem]{Fact}

\theoremstyle{definition}
\newtheorem{definition}[theorem]{Definition}
\newtheorem{example}[theorem]{Example}
\newtheorem{observation}[theorem]{Observation}
\newtheorem{assumption}[theorem]{Assumption}
\newtheorem*{assumption*}{Assumption}
\newtheorem{condition}[theorem]{Condition}

\theoremstyle{remark}
\newtheorem{remark}[theorem]{Remark}

\crefname{lemma}{Lemma}{Lemmas}
\crefname{theorem}{Theorem}{Theorems}
\crefname{definition}{Definition}{Definitions}
\crefname{fact}{Fact}{Facts}
\crefname{claim}{Claim}{Claims}
\crefname{proposition}{Proposition}{Propositions}


\newcommand{\dif}{\,\mathrm{d}}
%\newcommand{\E}{\mathbb{E}}
%\newcommand{\Var}{\mathrm{Var}}
\newcommand{\Cov}{\mathrm{Cov}}
\newcommand{\MI}{\mathrm{I}}
%\newcommand{\Ent}{\mathrm{Ent}}
\newcommand{\ind}{\mathbbm{1}}
\newcommand{\allone}{\mathbf{1}}
\newcommand{\tr}{\mathrm{tr}}
\newcommand{\wt}{\mathrm{weight}}
\DeclareMathOperator*{\argmax}{arg\,max}
\newcommand{\norm}[1]{\left\lVert #1 \right\rVert}
\newcommand{\TV}[2]{\left\|#1 - #2\right\|_{\mathrm{TV}}}
\newcommand{\ceil}[1]{\left\lceil #1 \right\rceil}
\newcommand{\floor}[1]{\left\lfloor #1 \right\rfloor}
\newcommand{\ip}[2]{\left\langle #1 , #2 \right\rangle}
\newcommand{\grad}{\nabla}
\newcommand{\hessian}{\nabla^2}
\newcommand{\trans}{\intercal}
\newcommand{\diag}{\mathrm{diag}}
\newcommand{\poly}{\mathrm{poly}}
\newcommand{\dist}{\mathrm{dist}}
\newcommand{\sgn}{\mathrm{sgn}}
\newcommand{\fpras}{\mathsf{FPRAS}}
\newcommand{\fpaus}{\mathsf{FPAUS}}
\newcommand{\fptas}{\mathsf{FPTAS}}

\newcommand{\eps}{\varepsilon}

\newcommand{\N}{\mathbb{N}}
\newcommand{\Z}{\mathbb{Z}}
\newcommand{\Q}{\mathbb{Q}}
\newcommand{\R}{\mathbb{R}}
\newcommand{\C}{\mathbb{C}}

\renewcommand{\AA}{\mathcal{A}}
\newcommand{\BB}{\mathcal{B}}
\newcommand{\CC}{\mathcal{C}}
\newcommand{\DD}{\mathcal{D}}
\newcommand{\EE}{\mathcal{E}}
\newcommand{\FF}{\mathcal{F}}
\newcommand{\GG}{\mathcal{G}}
\newcommand{\HH}{\mathcal{H}}
\newcommand{\II}{\mathcal{I}}
\newcommand{\JJ}{\mathcal{J}}
\newcommand{\KK}{\mathcal{K}}
\newcommand{\LL}{\mathcal{L}}
\newcommand{\MM}{\mathcal{M}}
\newcommand{\NN}{\mathcal{N}}
\newcommand{\OO}{\mathcal{O}}
\newcommand{\PP}{\mathcal{P}}
\newcommand{\QQ}{\mathcal{Q}}
\newcommand{\RR}{\mathcal{R}}
\renewcommand{\SS}{\mathcal{S}}
\newcommand{\TT}{\mathcal{T}}
\newcommand{\UU}{\mathcal{U}}
\newcommand{\VV}{\mathcal{V}}
\newcommand{\WW}{\mathcal{W}}
\newcommand{\XX}{\mathcal{X}}
\newcommand{\YY}{\mathcal{Y}}
\newcommand{\ZZ}{\mathcal{Z}}

\newcommand{\KL}[2]{D_{\mathrm{KL}} \left( #1 \,\Vert\, #2 \right)}
\newcommand{\KLbig}[2]{D_{\mathrm{KL}} \left( #1 \,\big\Vert\, #2 \right)}
\newcommand{\KLBig}[2]{D_{\mathrm{KL}} \left( #1 \,\Big\Vert\, #2 \right)}

\newcommand{\SKL}[2]{D_{\mathrm{SKL}} \left( #1 \,\Vert\, #2 \right)}
\newcommand{\SKLbig}[2]{D_{\mathrm{SKL}} \left( #1 \,\big\Vert\, #2 \right)}
\newcommand{\SKLBig}[2]{D_{\mathrm{SKL}} \left( #1 \,\Big\Vert\, #2 \right)}

\newcommand{\qandq}{\quad\text{and}\quad}
\newcommand{\supp}{\mathrm{supp}}

\newcommand{\DKL}[2]{\-D_{\-{KL}}\left(#1 \parallel #2\right)}
\newcommand{\DTV}[2]{\-D_{\mathrm{TV}}\left({#1},{#2}\right)}
% \newcommand{\dist}{\mathrm{dist}}
\newcommand{\e}{\mathrm{e}}
\renewcommand{\epsilon}{\varepsilon}
\renewcommand{\emptyset}{\varnothing}
% \newcommand{\norm}[1]{\left\Vert#1\right\Vert}
\newcommand{\set}[1]{\left\{#1\right\}}
\newcommand{\tuple}[1]{\left(#1\right)} 
% \newcommand{\eps}{\varepsilon}
\newcommand{\inner}[2]{\left\langle #1,#2\right\rangle}
\newcommand{\tp}{\tuple}
\newcommand{\ol}{\overline}
\newcommand{\abs}[1]{\left\vert#1\right\vert}
\newcommand{\ctp}[1]{\left\lceil#1\right\rceil}
\newcommand{\ftp}[1]{\left\lfloor#1\right\rfloor}
\newcommand{\Ind}[1]{\-{Ind}_{#1}}
\newcommand{\todo}[1]{{\color{red} TODO: #1}}

\newcommand{\Cor}{\Psi^{\mathrm{sym}}}
\newcommand{\cor}{\Psi^{\-{cor}}}
\newcommand{\Inf}{\Psi}
% \newcommand{\diag}{\-{diag}}

\def\*#1{\boldsymbol{#1}} % Use \*A for \mathbf{A}
\def\+#1{\mathcal{#1}} % Use \+A for \mathcal{A}
\def\-#1{\mathrm{#1}} % Use \-A for \mathrm{A}
\def\=#1{\mathbb{#1}} % Use \^A for \mathbb{A}
\def\!#1{\mathfrak{#1}} % Use \!A for \mathfrak{A}

\newcommand*\diff{\mathop{}\!\mathrm{d}}

% probability
% \DeclareMathOperator*{\oPr}{\mathbf{Pr}}
\def\oPr{\mathop{\mathrm{Pr}}}
\renewcommand{\Pr}[2][]{ \ifthenelse{\isempty{#1}}
  {\oPr\left[#2\right]}
  {\oPr_{#1}\left[#2\right]} } % Use \Pr[a]{b} for \mathbf{Pr}_a[b], \Pr{b} for  \mathbf{Pr}[b]

% expectation: relies on the package "dsfont"
% \DeclareMathOperator*{\oE}{\mathbb{E}}
\def\oE{\mathop{\mathbb{E}}}
\newcommand{\E}[2][]{ \ifthenelse{\isempty{#1}}
  {\oE\left[#2\right]}
  {\oE_{#1}\left[#2\right]} }

% variance
%\DeclareMathOperator*{\oVar}{\mathbf{Var}}
\def\oVar{\mathrm{Var}}
\newcommand{\Var}[2][]{ \ifthenelse{\isempty{#1}}
  {\oVar\left[#2\right]}
  {\oVar_{#1}\left[#2\right]} }

% entropy
% \DeclareMathOperator*{\oEnt}{\mathbf{Ent}}
\def\oEnt{\mathrm{Ent}}
\newcommand{\Ent}[2][]{ \ifthenelse{\isempty{#1}}
  {\oEnt\left[#2\right]}
  {\oEnt_{#1}\left[#2\right]} }

\newcommand{\PhiEnt}[2][]{ \ifthenelse{\isempty{#1}}
  {\oEnt^\phi\left[#2\right]}
  {\oEnt^\phi_{#1}\left[#2\right]} }


\newenvironment{Exampleparagraph}[1][]
{
    \par 
    \noindent 
    \textbf{\color{blue} #1}
    \par 
    \vspace{0.5em} 
    \itshape 
}
{
    \par 
    \vspace{1em} 
}

%\input{local-macro}

\newcommand{\RED}[1]{\textcolor{red}{#1}}
\newcommand{\zc}[1]{\textcolor{red}{ZC: #1}}

\newcommand{\bern}[1]{\mathrm{Bern}\left(#1\right)}
\newcommand{\pll}{\parallel}



\title{Degree-free Holant SI Reading}
\date{\today}

\pgfplotsset{compat=1.18} 
\begin{document}

\maketitle

\section{Main structure}
Assumptions:
\begin{itemize}
    \item for any $v\in V$, the signature of $v$ is log-concave.
    \item for $e=uv\in E$, let $\theta_{e}(i,j) = \lambda_e \frac{f_u(i+1)}{f_u(i)}\frac{f_v(j+1)}{f_v(j)}$ be the increased weight if we add the edge $e$. We have that $\theta_e(i,j)\in [\gamma,\Theta]$, where $\gamma$ and $\Theta$ are some constants independent of the graph.
    % \item 
\end{itemize}
Then we have the upper bound of marginal distribution:
\begin{equation}
    \Pr[S\sim \mu]{e\in S} = p_e \leq \frac{\Theta}{1+\Theta} .
\end{equation}

Let us fix $a\in \R^E$.
Define $\phi_e(i,j)$ backwards on $(\{0\}\cup [c_u])\times (\{0\}\cup [c_v])$ as follows,
\begin{align}
    \phi_e(i,j) = \begin{cases}
        0 &, i=c_u \lor j=c_v
        \\\theta_e(i,j)(1 - \phi_e(i+1,j+1)) &, \text{o.w.}
    \end{cases} ,
\end{align}
and $K_{e} := \phi_e(d_S(u),d_S(v))$ for $uv=e\in E$.
Here $\phi$ is \emph{intentionally} defined to have the following property:
\begin{align}\label{eq:1}
    \E[\mu]{K_e\mid X_{-e}} = \E[\mu]{X_e\mid X_{-e}} \implies \E[\mu]{K_eX_f} = \E[\mu]{X_eX_f }, \forall f\neq e .
\end{align}

Define 
$$
    F_a := \sum_{e\in E}a_eX_e\quad \text{ and }\quad K_a := \sum_{e\in E}a_eK_e.
$$
By \Cref{eq:1}, we decompose $\Var[\mu]{F_a}$ as follows,
\begin{align*}
    \Var[\mu]{F_a} &= \Cov(F_a,K_a) + \E[\mu]{F_a^2} - \E[\mu]{F_aK_a}
    \\&=\Cov(F_a,K_a) + \sum_{e_1,e_2\in E}a_{e_1}a_{e_2}(\E[\mu]{X_{e_1}X_{e_2}} - \E[\mu]{X_{e_1}K_{e_2}})
    \\&=\Cov(F_a,K_a) + \sum_{e\in E}a_{e}^2(\E[\mu]{X_e^2} - \E[\mu]{X_{e}K_{e}})
    \\&\leq \sqrt{\Var[\mu]{F_a}\Var[\mu]{K_a}} + \sum_{e\in E}a_{e}^2(p_e - \E[\mu]{X_{e}K_{e}}).
\end{align*}
Here the off-diagonal elements cancel and we only need to focus on the diagonal elements.
\begin{proposition}
    There exists some constant $H = H(b,\gamma,\Theta,\lambda)$ such that $\abs{K_e}\leq H$ for $e\in E$.
\end{proposition}
Define
$$
    Q(a) := \sum_{e\in E}a_e^2 p_e = a^\top \diag\{p\} a_e.
$$
we have that $\E[\mu]{-X_e K_e}\leq \max_{e\in E}\abs{K_e}\Pr[\mu]{e\in S} \leq p_eH$.
Therefore, by AM-GM inequality,
\begin{equation}
    \Var[\mu]{F_a} \leq \Var[\mu]{K_a} + (1+H)Q(a).
\end{equation}

\subsection{Upper bound $\Var[\mu]{K_a}$}

\begin{lemma}
    Consider two Holant measures that agree except at one vertex $v$. where their signatures are either
    $$
        (f_v, \+D f_v)\quad  \text{ or }\quad (\delta_{k+1},\delta_k),
    $$
    then \textbf{log-concavity} implies there is a coupling $(S^+,S^-)$ such that
    $$
        \sum_{x\in V} \abs{d_{S^+}(x) - d_{S^-}(x)} \leq 2.
    $$
\end{lemma}
Let $Y = \*1_{c_v - d_S(v)\geq r}$ where $c_v$ is the support length of $f_v$.
\begin{proposition}
    The above coupling implies the spectral stability of $Y$.
\end{proposition}
\begin{lemma}\label{lem:1}
    The spectral stability of $Y$ ($\Cov(Y)\preceq C\cdot \diag\{\E Y\}$) implies the spectral gap of some resulting kernel.
    Together with spectral stability, the spectral gap implies the following bound
    $$
        \Var[\mu]{K}\leq 4^b\tp{\sum_i\E[\mu]{Y_iH_i^2} + \sum_{ij}c_{ij}\E[\mu]{Y_iY_j}},
    $$
    where $K(Y) = \sum_{ij}c_{ij}Y_iY_j$ is the quadratic form of $Y$ and $H_i(Y) = \sum_j c_{ij}Y_j$ is the ``first derivative''.
\end{lemma}
We can also express $K_a$ as the quadratic form of $Y$ variables:
\begin{align}
    K_e = \sum_{e\in E}a_e\sum_{r=1}^{c_u}\sum_{s=1}^{c_v} \alpha_{rs}^e Y_{u,r} Y_{v,s},
\end{align}
where $\alpha$ is ``second derivative" of $\phi_e$.
Define 
$$
    R(a) = \sum_{v\in V} \frac{(\sum_{e:v\in e} p_ea_e)^2}{\E[\mu]{Y_{v,1}}} .
$$
Plugging \Cref{lem:1} into $K_e$, we have that
$$
    \Var[\mu]{K_e}\lesssim \underbrace{\sum_{v,s=1}\E[\mu]{Y_{v,s}H_{v,s}^2}}_{\lesssim Q(a)+R(a)} + \underbrace{\sum_{v,s\geq 2}\E[\mu]{Y_{v,s}H_{v,s}^2}}_{\lesssim Q(a)} + \underbrace{\sum_{u,v,s,r}c_{ij}\E[\mu]{Y_{u,s}Y_{v,r}}}_{\lesssim Q(a)}
$$
\color{red}
Use Hodge decomposition to bound $R(a)$.
\end{document}